%%=============================================================================
%% Samenvatting
%%=============================================================================

% TODO: De "abstract" of samenvatting is een kernachtige (~ 1 blz. voor een
% thesis) synthese van het document.
%
% Een goede abstract biedt een kernachtig antwoord op volgende vragen:
%
% 1. Waarover gaat de bachelorproef?
% 2. Waarom heb je er over geschreven?
% 3. Hoe heb je het onderzoek uitgevoerd?
% 4. Wat waren de resultaten? Wat blijkt uit je onderzoek?
% 5. Wat betekenen je resultaten? Wat is de relevantie voor het werkveld?
%
% Daarom bestaat een abstract uit volgende componenten:
%
% - inleiding + kaderen thema
% - probleemstelling
% - (centrale) onderzoeksvraag
% - onderzoeksdoelstelling
% - methodologie
% - resultaten (beperk tot de belangrijkste, relevant voor de onderzoeksvraag)
% - conclusies, aanbevelingen, beperkingen
%
% LET OP! Een samenvatting is GEEN voorwoord!

%%---------- Nederlandse samenvatting -----------------------------------------
%
% TODO: Als je je bachelorproef in het Engels schrijft, moet je eerst een
% Nederlandse samenvatting invoegen. Haal daarvoor onderstaande code uit
% commentaar.
% Wie zijn bachelorproef in het Nederlands schrijft, kan dit negeren, de inhoud
% wordt niet in het document ingevoegd.

\IfLanguageName{english}{%
\selectlanguage{dutch}
\chapter*{Samenvatting}
\lipsum[1-4]
\selectlanguage{english}
}{}

%%---------- Samenvatting -----------------------------------------------------
% De samenvatting in de hoofdtaal van het document

\chapter*{\IfLanguageName{dutch}{Samenvatting}{Abstract}}

Bedrijven die afhankelijk zijn van fysieke locaties hebben vaak onvoldoende inzicht in de aantrekkelijkheid van een bepaalde locatie voor huidige en potentiële klanten. Hierdoor lopen ondernemingen het risico verkeerde beslissingen te nemen bij het openen van nieuwe vestigingen of het inkopen van goederen. Dit kan leiden tot lage bezoekersaantallen of financiële schade. Dit onderzoek toont aan in welke mate artificiële intelligentie (AI) op basis van censusdata de populariteit van een point-of-interest (POI) betrouwbaar kan voorspellen. Op die manier kunnen bedrijven hun vestigingen en inkopen gerichter afstemmen op veelbelovende locaties en het verwachte klantenaantal, waardoor hun marketingactiviteiten effectiever en efficiënter worden.

Om deze onderzoeksvraag te beantwoorden wordt nagegaan welke neurale netwerken geschikt zijn om de populariteit van een POI te voorspellen op basis van census data. Hierbij wordt gedefinieerd wat censusdata inhoudt en in welke mate dit een verklarende variabele is voor bezoekersaantallen. Verder worden onderzoekspapers bestudeerd om de voor- en nadelen van deze neurale netwerken te achterhalen. Daarnaast wordt er onderzocht welke methoden ingezet kunnen worden om de prestaties van modellen te optimaliseren. Na de literatuurstudie is een proof of concept (POC) ontwikkeld. In deze POC worden neurale netwerken ontwikkeld, getraind en getest om de populariteit van een POI mee te voorspellen. Hierbij worden de bevindingen uit de literatuurstudie toegepast.

De resultaten uit de POC suggereren dat de populariteit van een POI op zekere mate voorspeld kan worden met de huidige bestaande neurale netwerken. In deze POC zijn MLP, RNN en Transformer modellen getraind en geëvalueerd, waarbij de Transformer het beste presteert op de gecreëerde dataset. Hierbij levert de Transformer de beste prestaties op de diverse experimenten en presteert het model ook het beste op verschillende termijnen. Hoewel er nog beperkingen zijn en ruimte voor verbetering bestaat, bevestigt deze studie de meerwaarde om de populariteit van een POI te voorspellen aan de hand van neurale netwerken. De bevindingen hebben invloed op alle aspecten van de bedrijfswereld. Door de populariteit van een locatie te voorspellen kunnen bedrijven strategische beslissingen nemen bij een vestiging te openen en gerichter goederen inkopen. Hierdoor zullen ondernemingen meer omzet kunnen maken en startende bedrijven zullen een grotere kans hebben om te overleven. 
